\documentclass[11pt,a4paper]{article}

\usepackage[T1]{fontenc}
\usepackage{lmodern}
\usepackage{amsmath,amssymb}
\usepackage[hmargin=0.9in,vmargin=0.75in]{geometry}
\usepackage[hidelinks]{hyperref}

\setlength{\parindent}{0pt}
\setlength{\parskip}{0.4em}
\DeclareMathOperator*{\argmin}{arg\,min}
\DeclareMathOperator*{\argmax}{arg\,max}

\title{Why Gaussian Noise Leads to Least Squares}
\author{}
\date{}

\begin{document}
\maketitle
\vspace{-2em}

\textbf{The short answer:} the Gaussian probability density contains a
\emph{squared error} in its exponent. Independent observations give a
\emph{product} of densities. Taking the logarithm turns that product into a
\emph{sum}, so maximizing likelihood becomes minimizing a sum of squared errors.

\section{Start with a model for the errors}

Suppose we observe pairs $(x_i,y_i)$, for $i=1,\ldots,n$, and predict $y_i$ using
$f_\theta(x_i)$, where $\theta$ denotes the parameters we want to fit. Assume
\[
    y_i=f_\theta(x_i)+\varepsilon_i,
    \qquad \varepsilon_i\overset{\mathrm{iid}}{\sim}\mathcal{N}(0,\sigma^2).
\]
Here ``iid'' means independent and identically distributed: the errors are
independent, have mean zero, and share the same variance $\sigma^2>0$.
Treat the inputs $x_i$ as fixed and, for now, treat $\sigma^2$ as fixed too.
The Gaussian probability density function (PDF) for one observation is
\begin{equation}
    p(y_i\mid x_i,\theta,\sigma^2)
    =\frac{1}{\sqrt{2\pi\sigma^2}}
      \exp\!\left[-\frac{\bigl(y_i-f_\theta(x_i)\bigr)^2}{2\sigma^2}\right].
\end{equation}
The \textbf{residual} is $r_i(\theta)=y_i-f_\theta(x_i)$: observed minus predicted.
The square is already present in the Gaussian PDF; taking logs does not create it.

\section{Multiply the densities, then take the logarithm}

The \textbf{likelihood} is the joint density of the observed data, viewed as a
function of the parameters. Independence lets us multiply the individual densities:
\begin{align}
    L(\theta)
    &=\prod_{i=1}^{n}p(y_i\mid x_i,\theta,\sigma^2)\\
    &=(2\pi\sigma^2)^{-n/2}
      \exp\!\left[-\frac{1}{2\sigma^2}
                    \sum_{i=1}^{n}r_i(\theta)^2\right].
\end{align}
For continuous data, this is a density, not the probability of observing an
exact value (which is zero). Maximum likelihood compares these densities.

Because $\log$ is strictly increasing, maximizing $L$ is equivalent to maximizing
its log, $\ell(\theta)=\log L(\theta)$. Equivalently, minimize the negative log:
\begin{equation}
    -\ell(\theta)
    =\underbrace{\frac{n}{2}\log(2\pi\sigma^2)}_{\text{constant with respect to }\theta}
     +\underbrace{\frac{1}{2\sigma^2}}_{\text{positive constant}}
      \underbrace{\sum_{i=1}^{n}\bigl(y_i-f_\theta(x_i)\bigr)^2}_{\displaystyle S(\theta):\ \text{sum of squared errors}}.
\end{equation}
Adding a constant or multiplying by a positive constant does not change the
minimizer. Therefore,
\begin{equation}
    \boxed{\widehat\theta_{\mathrm{MLE}}
    \in\argmax_\theta L(\theta)
    =\argmin_\theta\sum_{i=1}^{n}\bigl(y_i-f_\theta(x_i)\bigr)^2.}
\end{equation}
This is \textbf{least squares}. Dividing the sum by $n$ gives mean squared error
(MSE), with exactly the same minimizers.

\newpage
\section{A concrete example: estimating a constant mean}

If $y_i\overset{\mathrm{iid}}{\sim}\mathcal{N}(\mu,\sigma^2)$, the prediction is simply
$f_\mu(x_i)=\mu$. Least squares minimizes
\[
    S(\mu)=\sum_{i=1}^{n}(y_i-\mu)^2.
\]
Differentiate and set the derivative to zero:
\[
    \frac{dS}{d\mu}=2\sum_{i=1}^{n}(\mu-y_i)=0
    \quad\Longrightarrow\quad
    \boxed{\widehat\mu=\frac{1}{n}\sum_{i=1}^{n}y_i.}
\]
Since $d^2S/d\mu^2=2n>0$, this is the unique minimum. Thus the sample mean is
also the Gaussian maximum-likelihood estimate of the population mean.
For observations $1,2,3$, fitting $\mu=2$ gives $S=2$, whereas $\mu=1$ gives $S=5$.
Their likelihood ratio is
\[
    \frac{L(2)}{L(1)}
    =\exp\!\left(\frac{5-2}{2\sigma^2}\right)>1:
\]
the smaller squared error means the larger likelihood.
For a line, replace $\mu$ by $a x_i+b$; the same derivation yields the usual
linear-regression objective $\sum_i(y_i-a x_i-b)^2$.

\section{What if the variance is unknown?}

If we estimate a single shared variance as well, we must retain the terms
involving $\sigma^2$. Writing $v=\sigma^2$, the negative log-likelihood is
\[
    J(\theta,v)=\frac{n}{2}\log(2\pi v)+\frac{S(\theta)}{2v}.
\]
For a fixed $\theta$ with $S(\theta)>0$, minimizing over $v>0$ gives
\[
    \frac{\partial J}{\partial v}
    =\frac{n}{2v}-\frac{S(\theta)}{2v^2}=0
    \quad\Longrightarrow\quad
    \widehat v(\theta)=\frac{S(\theta)}{n}.
\]
Substituting back yields
\[
    J\bigl(\theta,\widehat v(\theta)\bigr)
    =\frac{n}{2}\left[\log\!\left(2\pi\frac{S(\theta)}{n}\right)+1\right].
\]
This is increasing in $S(\theta)$, so minimizing it still gives least squares.
When the minimum squared error is positive and attained,
$\widehat\sigma^2_{\mathrm{MLE}}=S(\widehat\theta)/n$.
The denominator here is $n$, not a degrees-of-freedom correction used for
unbiased variance estimation. If a perfect fit gives $S=0$, the likelihood is
unbounded as $\sigma^2\to0$; there is no finite joint MLE with positive variance.

\section{Which assumptions make ordinary least squares appear?}

\begin{itemize}
    \item \textbf{Gaussian errors} supply the squared residual. Independent
    Laplace errors with a fixed common scale instead give absolute-error loss,
    $\sum_i |r_i|$.
    \item \textbf{Independence} allows the likelihood to factor into a product,
    so its log becomes a sum of individual contributions.
    \item \textbf{Equal variance} makes every squared residual have the same
    weight. Independent Gaussian errors with known unequal variances
    $\sigma_i^2$ instead give weighted least squares,
    $\sum_i r_i^2/\sigma_i^2$.
\end{itemize}

\textbf{Takeaway:} least squares can be used without assuming Gaussian noise,
but independent Gaussian errors with a common variance give it a precise
maximum-likelihood interpretation.

\end{document}
